\chapter*{ABSTRACT}
\addcontentsline{toc}{chapter}{ABSTRACT}

\begin{otherlanguage}{english}
This project designed a machine to clean large bearings from high-pressure grinding rolls (HPGR) in mining maintenance workshops. The reference was a bearing of 3032 kg, 1420 mm outer diameter and 515 mm width. Following the VDI 2221 product development model, fourteen specifications were defined and four alternatives were generated. A decision matrix selected a configuration in which the bearing remains stationary on a three-support table, moved by a double chain, while a rotating U-shaped distributor and an internal lance deliver the fluid through 24 nozzles. The support was verified with a frame structural analysis, a global model including non-uniform torsion, finite element models of the centering units and the design of its welded joints. The hydraulic network was solved by successive approximations, and wetting was simulated with a discrete phase model verified on three meshes. The frame reached a maximum demand-to-capacity ratio of 0.82, and the bearing placement error was 4.24 mm against a 5 mm limit. The network required 221 L/min at about 53 m, with nozzle pressures from 4.00 to 4.13 bar. A pair of narrow-angle nozzles increased the dose on the outer zone of the lower face, the critical surface, by a factor of 4.8 and reduced the estimated washing time of that face, which sets the machine time, from 37.6 to 14.1 min. The investment budget was estimated at PEN 314,310, excluding VAT. The design is feasible as preliminary engineering verified by simulation and meets thirteen of the fourteen specifications in the models; its cleaning effectiveness and economic justification require tests and data beyond the scope of the project.

\medskip
\noindent\textbf{Keywords:} \textit{HPGR bearing; spray cleaning; machine design; computational fluid dynamics; finite elements; mining maintenance.}
\end{otherlanguage}
